% =================================================================
% PREAMBULE MATH LYCÉE — ARCHIVÉ
% Ce fichier n'est plus chargé dans les cours de collège.
% À réutiliser si nécessaire pour des cours de lycée.
% =================================================================

% --- Ensembles (déjà dans Preambule_Commun.tex pour le collège) ---
\newcommand{\N}{\ensuremath{\mathbf{N}}}
\newcommand{\Z}{\ensuremath{\mathbf{Z}}}
\newcommand{\D}{\ensuremath{\mathbf{D}}}
\newcommand{\Q}{\ensuremath{\mathbf{Q}}}
\newcommand{\R}{\ensuremath{\mathbf{R}}}
\newcommand{\card}{\text{card}}

\newcommand{\eq}{\Longleftrightarrow}
\newcommand{\abs}[1]{\left\lvert #1 \right\rvert}
\renewcommand{\vec}{\overrightarrow}
\newcommand{\norme}[1]{\left\lVert \vec{#1} \right\rVert}
\DeclareMathOperator{\e}{e}

% --- Vecteurs et coordonnées ---
\newcommand{\ivec}{\vec{\imath}}
\newcommand{\jvec}{\vec{\jmath}}
\newcommand{\Oij}{\ensuremath{(O, \ivec, \jvec)}}
\newcommand{\veccol}[3]{\vec{#1} = \begin{pmatrix} #2 \\ #3 \end{pmatrix}}
\newcommand{\veccoo}[2]{\begin{pmatrix} #1 \\ #2 \end{pmatrix}}

% --- Intervalles (TikZ) ---
\usepackage{ifthen}

% Borne inf; ouverte/fermée (o/f); borne sup; ouverte/fermée (o/f)
\newcommand{\intervalleborne}[4]{%
\begin{tikzpicture}[baseline=-0.5ex]
    \draw[->] (0,0) -- (5,0);
    \draw[line width=1mm, green!50!black] (1,0) -- (4,0);
    \node[below=5pt] at (1,0) {#1};
    \node[below=5pt] at (4,0) {#3};
    \node[green!50!black] at (1,0)
        {\ifthenelse{\equal{\detokenize{#2}}{\detokenize{f}}}{$\bigg[$}{$\bigg]$}};
    \node[green!50!black] at (4,0)
        {\ifthenelse{\equal{\detokenize{#4}}{\detokenize{f}}}{$\bigg]$}{$\bigg[$}};
\end{tikzpicture}}

% Côté ouvert (-/+); borne réelle; ouverte/fermée (o/f)
\newcommand{\intervallenonborne}[3]{%
\begin{tikzpicture}[baseline=-0.5ex]
    \draw[->] (0,0) -- (5,0);
    \ifthenelse{\equal{\detokenize{#1}}{\detokenize{-}}}{%
        \draw[line width=1mm, green!50!black] (0,0) -- (4,0);
        \node[below=5pt] at (0,0) {$-\infty$};
        \node[below=5pt] at (4,0) {#2};
        \node[green!50!black] at (4,0)
            {\ifthenelse{\equal{\detokenize{#3}}{\detokenize{f}}}{$\bigg]$}{$\bigg[$}};
    }{%
        \draw[line width=1mm, green!50!black] (1,0) -- (5,0);
        \node[below=5pt] at (1,0) {#2};
        \node[below=5pt] at (5,0) {$+\infty$};
        \node[green!50!black] at (1,0)
            {\ifthenelse{\equal{\detokenize{#3}}{\detokenize{f}}}{$\bigg[$}{$\bigg]$}};
    }
\end{tikzpicture}}

% --- Segments et diagrammes ---
\newcommand{\segmentlegende}[4]{%
\begin{tikzpicture}[baseline=-0.5ex]
    \draw[->] (0,0) -- (5,0);
    \draw (1,-0.1) -- (1,0.1) node[above=5pt] {#1} node[below=5pt] {#2};
    \draw (4,-0.1) -- (4,0.1) node[above=5pt] {#3} node[below=5pt] {#4};
    \draw[<->] (1,0.8) -- node[above] {$AM = #4 - #2$} (4,0.8);
\end{tikzpicture}}

\newcommand{\diagrammedouble}[6]{%
\begin{tikzpicture}[baseline=(B.base)]
    \node(A) at (0,0) {#1}; \node(B) at (4,0) {#2}; \node(C) at (8,0) {#3};
    \draw[->, dashed, line width=0.8mm, green!50!black]
        (A) to[bend left=20] node[midway, above] {#4} (B);
    \draw[->, dashed, line width=0.8mm, blue!50!black]
        (B) to[bend left=20] node[midway, above] {#5} (C);
    \draw[->, dashed, line width=0.8mm, red!50!black]
        (A) to[bend right=25] node[midway, below] {#6} (C);
\end{tikzpicture}}
